\section{Đánh giá các Tính năng AI}

\subsection{Tại sao Kiểm thử Truyền thống Không Đủ cho AI}

Phương pháp luận kiểm thử được mô tả trong chương trước tuân theo mô hình \textbf{định tính (deterministic)}: với bất kỳ đầu vào cho trước nào, hệ thống phải tạo ra đầu ra được xác định chính xác, và một phép kiểm tra (assertion) chỉ có thể đạt (pass) hoặc thất bại (fail). Mô hình này hoạt động tốt cho các hàm không lưu trạng thái (stateless) và các điểm cuối REST API nơi hành vi mang tính dự đoán được và có giới hạn.

Tuy nhiên, các tính năng tích hợp AI trong BkLib — cụ thể là Chatbot dựa trên RAG và Động cơ Tìm kiếm Ngữ nghĩa — tạo ra các đầu ra \textbf{không định tính (non-deterministic)} dưới dạng ngôn ngữ tự nhiên. Hãy xem xét sự đối lập sau:

\begin{table}[H]
\centering
\caption{Kiểm thử Định tính vs. Đánh giá Tính năng AI}
\begin{tabularx}{\textwidth}{|p{3.5cm}|X|X|}
\hline
\textbf{Khía cạnh} & \textbf{Kiểm thử Truyền thống} & \textbf{Đánh giá Tính năng AI} \\
\hline
Kiểu đầu ra & Giá trị cố định (mã trạng thái, trường JSON) & Câu ngôn ngữ tự nhiên hoặc danh sách xếp hạng \\
\hline
Kiểm tra tính đúng đắn & \texttt{expect(x).toBe(y)} & Độ tương đồng ngữ nghĩa, độ trung thực, độ liên quan \\
\hline
Mức độ biến đổi & Không có — đầu vào giống nhau cho đầu ra giống nhau & Cao — diễn đạt thay đổi qua từng lần chạy \\
\hline
Công cụ & Jest, Supertest & LLM-as-a-Judge, độ tương đồng cosine vector \\
\hline
\end{tabularx}
\end{table}

\noindent Ví dụ, khi độc giả hỏi \textit{``Tôi có thể mượn tối đa bao nhiêu quyển sách cùng lúc?''}, phản hồi hợp lệ của chatbot có thể là \textit{``Bạn có thể mượn tối đa 5 quyển sách cho mỗi lần mượn''} hoặc \textit{``Thư viện quy định hạn mức tối đa 5 tài liệu cho mỗi độc giả''} — cả hai đều đúng về mặt ngữ nghĩa nhưng khác nhau về mặt từ ngữ. Phép kiểm tra cứng \texttt{toBe()} không thể nắm bắt được sự tương đương này, do đó cần phải có một khung đánh giá AI chuyên biệt.

\subsection{Phương pháp luận AI-as-a-Judge}

Đồ án áp dụng phương pháp luận \textbf{LLM-as-a-Judge} (sử dụng một LLM mạnh hơn để đánh giá đầu ra của hệ thống) kết hợp với truy vấn vector trực tiếp để đánh giá định lượng cả hai dịch vụ AI.

%\begin{figure}[H]
%\centering
%\includegraphics[width=0.85\textwidth]{Images/ai_eval_architecture.png}
%\caption{Kiến trúc Đánh giá AI}
%\label{fig:ai_eval_arch}
%\end{figure}

\subsubsection{Các Chỉ số Đánh giá Chatbot (Khung RAG)}

Mô-đun Chatbot sử dụng RAG được đánh giá trên ba chỉ số cốt lõi, mỗi chỉ số được chấm trên thang điểm 1--5 bởi một mô hình giám khảo độc lập:

\begin{enumerate}
    \item \textbf{Độ trung thực (Faithfulness) [1--5]:} Đo lường xem câu trả lời của Chatbot có hoàn toàn dựa trên các đoạn văn bản ngữ cảnh được truy xuất hay không, hay chứa thông tin suy diễn không có thực (hallucination).
    \item \textbf{Độ liên quan (Answer Relevancy) [1--5]:} Đo lường mức độ câu trả lời giải quyết trực tiếp câu hỏi của người dùng mà không lan man hoặc trả lời lạc đề.
    \item \textbf{Độ chính xác (Correctness) [1--5]:} So sánh câu trả lời được tạo với câu trả lời chuẩn (Ground Truth) được gắn nhãn bởi con người về mặt ngữ nghĩa.
\end{enumerate}

\subsubsection{Các Chỉ số Đánh giá Tìm kiếm Ngữ nghĩa}

Tính năng Tìm kiếm Ngữ nghĩa được đánh giá dựa trên tiêu chuẩn truy xuất thông tin:

\begin{enumerate}
    \item \textbf{Top-K Accuracy (Hit@K):} Tỷ lệ phần trăm các câu truy vấn mà ít nhất một tài liệu liên quan xuất hiện trong $K$ kết quả đầu tiên (với $K=3$ và $K=5$).
    \item \textbf{Mean Reciprocal Rank (MRR):} Điểm trung bình của vị trí nghịch đảo cho tài liệu liên quan đầu tiên tìm thấy:
    \begin{equation}
    \text{MRR} = \frac{1}{|Q|} \sum_{i=1}^{|Q|} \frac{1}{\text{rank}_i}
    \end{equation}
\end{enumerate}

\subsection{Kết quả Đánh giá Chatbot (RAG)}

\subsubsection{Thiết lập Bộ dữ liệu Kiểm thử}

Bộ dữ liệu kiểm thử Chatbot bao gồm **16 ca kiểm thử** đại diện cho 4 nhóm nghiệp vụ chính:
\begin{itemize}
    \item \textbf{Quy định mượn trả (Borrowing Rules):} 5 ca kiểm thử (CB-01 đến CB-05).
    \item \textbf{Thủ tục gia hạn và đặt chỗ (Renewal \& Holds):} 4 ca kiểm thử (CB-06 đến CB-09).
    \item \textbf{Phí phạt quá hạn và chính sách (Fines \& Policy):} 4 ca kiểm thử (CB-10 đến CB-13).
    \item \textbf{Truy vấn ngoài phạm vi / Khóa tài khoản (Out-of-Scope / Account Lock):} 3 ca kiểm thử (CB-14 đến CB-16).
\end{itemize}

\subsubsection{Bảng Kết quả Đánh giá Chatbot Chi tiết}

Bảng~\ref{tab:chatbot_eval_results} trình bày chi tiết kết quả đánh giá cho toàn bộ 16 ca kiểm thử.

\begin{table}[H]
\centering
\caption{Kết quả Đánh giá Chi tiết Chatbot RAG (Thang điểm 1--5)}
\label{tab:chatbot_eval_results}
\begin{tabularx}{\textwidth}{|l|p{3.2cm}|X|c|c|c|}
\hline
\textbf{ID} & \textbf{Nhóm Nghiệp vụ} & \textbf{Nội dung Câu hỏi} & \textbf{Faith.} & \textbf{Relev.} & \textbf{Correct.} \\
\hline
CB-01 & Mượn trả & Mượn tối đa bao nhiêu quyển sách? & 5 & 5 & 5 \\
\hline
CB-02 & Mượn trả & Sinh viên được mượn sách trong bao lâu? & 5 & 5 & 5 \\
\hline
CB-03 & Mượn trả & Giảng viên có hạn mức mượn khác không? & 5 & 5 & 5 \\
\hline
CB-04 & Mượn trả & Tôi có thể mượn tạp chí in không? & 5 & 5 & 5 \\
\hline
CB-05 & Mượn trả & Quy trình mượn sách trực tiếp tại quầy? & 5 & 4 & 5 \\
\hline
CB-06 & Gia hạn \& Đặt chỗ & Cách gia hạn sách trực tuyến? & 5 & 5 & 5 \\
\hline
CB-07 & Gia hạn \& Đặt chỗ & Sách có người đặt chỗ có gia hạn được không? & 5 & 5 & 5 \\
\hline
CB-08 & Gia hạn \& Đặt chỗ & Được gia hạn tối đa bao nhiêu lần? & 5 & 5 & 5 \\
\hline
CB-09 & Gia hạn \& Đặt chỗ & Làm sao biết sách đặt chỗ đã về? & 5 & 5 & 5 \\
\hline
CB-10 & Phí phạt & Mức phạt trả sách quá hạn 1 ngày? & 5 & 5 & 5 \\
\hline
CB-11 & Phí phạt & Ngày nghỉ lễ có tính phí phạt không? & 5 & 5 & 5 \\
\hline
CB-12 & Phí phạt & Thanh toán tiền phạt ở đâu? & 5 & 5 & 5 \\
\hline
CB-13 & Phí phạt & Bị phạt bao nhiêu thì bị khóa thẻ mượn? & 5 & 4 & 4 \\
\hline
CB-14 & Ngoài phạm vi & Cửa hàng canteen trường nằm ở đâu? & 5 & 5 & 5 \\
\hline
CB-15 & Ngoài phạm vi & Mật khẩu wifi thư viện là gì? & 5 & 5 & 5 \\
\hline
CB-16 & Ngoài phạm vi & Hướng dẫn giải bài tập môn Giải tích 1? & 4 & 4 & 4 \\
\hline
\multicolumn{3}{|r|}{\textbf{Trung bình Tổng thể}} & \textbf{4.94} & \textbf{4.88} & \textbf{4.88} \\
\hline
\end{tabularx}
\end{table}

\subsubsection{Phân tích Kết quả Đánh giá Chatbot}

%\begin{figure}[H]
%\centering
%\includegraphics[width=0.85\textwidth]{Images/ai_eval_scores.png}
%\caption{Phân bố Điểm số Đánh giá Chatbot RAG}
%\label{fig:chatbot_scores}
%\end{figure}

\begin{itemize}
    \item \textbf{Độ trung thực Trung bình (4.94/5):} Đạt điểm gần như tuyệt đối. Nhờ cơ chế RAG giới hạn ngữ cảnh nghiêm ngặt, Chatbot không gặp phải hiện tượng bịa đặt thông tin (hallucination) trong các câu hỏi nghiệp vụ.
    \item \textbf{Độ liên quan Trung bình (4.88/5):} Các phản hồi tập trung trực tiếp vào câu hỏi của độc giả, đưa ra câu trả lời ngắn gọn và hướng dẫn hành động rõ ràng.
    \item \textbf{Độ chính xác Trung bình (4.88/5):} So sánh với câu trả lời chuẩn (Ground Truth) xác nhận Chatbot truyền tải chính xác các quy định về hạn mức mượn, mức phạt và quy trình đặt chỗ.
\end{itemize}

\subsection{Kết quả Đánh giá Tìm kiếm Ngữ nghĩa (Semantic Search)}

\subsubsection{Thiết lập Bộ dữ liệu Kiểm thử}

Bộ dữ liệu kiểm thử Tìm kiếm Ngữ nghĩa bao gồm \textbf{15 ca kiểm thử} (12 câu truy vấn ngữ nghĩa có ý nghĩa + 3 câu truy vấn nhiễu/ngoài phạm vi), được thiết kế nhằm mô phỏng thói quen tìm kiếm thực tế của độc giả trong thư viện đại học. Các truy vấn được chia thành 4 nhóm:

\begin{itemize}
    \item \textbf{Nhóm 1 --- Kỹ thuật / Lập trình (S01--S05):} Truy vấn về Python, Học máy, thiết kế UI/UX, hệ thống phân tán và giải thuật.
    \item \textbf{Nhóm 2 --- Khoa học Xã hội / Kinh tế (S06--S08):} Truy vấn về quản lý dự án, kinh tế vĩ mô và tâm lý học.
    \item \textbf{Nhóm 3 --- Ngữ cảnh / Ngôn ngữ tự nhiên (S09--S12):} Các truy vấn tự nhiên như ``tôi muốn học cách viết code sạch'' hoặc ``sách lịch sử Việt Nam''.
    \item \textbf{Nhóm 4 --- Câu hỏi nhiễu (N01--N03):} Các đầu vào hoàn toàn lạc đề (thời tiết, chuỗi vô nghĩa, đặt đồ ăn) --- kỳ vọng hệ thống không trả về kết quả liên quan.
\end{itemize}

Mỗi truy vấn được đánh giá trên tập dữ liệu \textbf{72 cuốn sách} đã được index bằng pgvector, sử dụng ba chỉ số truy xuất (NDCG@5, MRR, Precision@3) và điểm đánh giá độ liên quan bởi LLM-as-a-Judge (thang 1--5).

\subsubsection{Tổng hợp Chỉ số Đánh giá}

\begin{table}[H]
\centering
\caption{Tổng hợp Chỉ số Đánh giá Tìm kiếm Ngữ nghĩa (pgvector)}
\label{tab:search_metrics_summary}
\begin{tabularx}{\textwidth}{|p{5.5cm}|X|c|c|}
\hline
\textbf{Chỉ số} & \textbf{Ý nghĩa} & \textbf{Kết quả} & \textbf{Mục tiêu} \\
\hline
NDCG@5 & Chất lượng xếp hạng top 5 kết quả & \textbf{0.9288} & $\geq 0.70$ \checkmark \\
\hline
MRR (Mean Reciprocal Rank) & Vị trí của kết quả đúng đầu tiên & \textbf{0.8750} & $\geq 0.60$ \checkmark \\
\hline
Precision@3 & Tỷ lệ kết quả đúng trong top 3 & \textbf{0.5556} & $\geq 0.60$ \\
\hline
Noise Rejection Rate & Tỷ lệ từ chối truy vấn không liên quan & \textbf{100.0\%} & $\geq 80\%$ \checkmark \\
\hline
LLM Relevance Score (Trung bình) & Điểm đánh giá ngữ nghĩa bởi LLM & \textbf{4.13 / 5} & --- \\
\hline
LLM Ranking Score (Trung bình) & Điểm đánh giá xếp hạng bởi LLM & \textbf{3.80 / 5} & --- \\
\hline
\end{tabularx}
\end{table}

\subsubsection{Bảng Kết quả Đánh giá Tìm kiếm Chi tiết}

Bảng~\ref{tab:search_eval_results} trình bày chi tiết kết quả cho toàn bộ 15 ca kiểm thử. Cột "Rel." và "Rank." là điểm số do mô hình LLM (Gemini) tự động chấm độc lập cho chất lượng Ngữ nghĩa (Relevance) và Xếp hạng (Ranking) của kết quả trả về, trên thang điểm 5.

\begin{table}[H]
\centering
\caption{Kết quả Đánh giá Chi tiết Tìm kiếm Ngữ nghĩa (pgvector, 15 truy vấn)}
\label{tab:search_eval_results}
\begin{tabularx}{\textwidth}{|l|X|c|c|c|c|c|c|c|}
\hline
\textbf{ID} & \textbf{Truy vấn} & \textbf{SL} & \textbf{Top\%} & \textbf{NDCG} & \textbf{MRR} & \textbf{P@3} & \textbf{Rel.} & \textbf{Rank.} \\
\hline
S01 & sách lập trình Python cho người mới... & 6 & 76.9 & 1.000 & 1.00 & 1.000 & 5 & 5 \\
\hline
S02 & tài liệu học máy học và trí tuệ... & 8 & 68.1 & 1.000 & 1.00 & 0.667 & 4 & 4 \\
\hline
S03 & sách về thiết kế giao diện và... & 8 & 61.4 & 1.000 & 0.00 & 0.000 & 1 & 1 \\
\hline
S04 & tài liệu hệ thống phân tán và... & 3 & 80.7 & 1.000 & 1.00 & 0.667 & 5 & 5 \\
\hline
S05 & sách cấu trúc dữ liệu và giải... & 8 & 72.1 & 1.000 & 1.00 & 0.333 & 4 & 5 \\
\hline
S06 & quản lý dự án công nghệ phần mềm & 8 & 71.1 & 0.920 & 1.00 & 0.667 & 4 & 3 \\
\hline
S07 & kinh tế vĩ mô và tài chính... & 8 & 71.5 & 0.833 & 1.00 & 0.667 & 4 & 3 \\
\hline
S08 & sách tâm lý học và kỹ năng... & 8 & 66.3 & 0.877 & 1.00 & 0.333 & 5 & 4 \\
\hline
S09 & muốn học cách viết code sạch & 1 & 67.9 & 1.000 & 1.00 & 0.333 & 5 & 5 \\
\hline
S10 & mạng máy tính và bảo mật & 7 & 69.2 & 0.920 & 1.00 & 0.667 & 4 & 3 \\
\hline
S11 & sách về cơ sở dữ liệu và SQL & 3 & 59.1 & 0.631 & 0.50 & 0.333 & 3 & 2 \\
\hline
S12 & sách lịch sử Việt Nam thời kỳ... & 8 & 85.4 & 0.965 & 1.00 & 1.000 & 3 & 2 \\
\hline
N01 & hôm nay thời tiết như thế nào & 0 & 27.7 & --- & --- & --- & 5 & 5 \\
\hline
N02 & abc xyz 123 & 0 & 32.5 & --- & --- & --- & 5 & 5 \\
\hline
N03 & tôi muốn đặt đồ ăn online & 0 & 22.6 & --- & --- & --- & 5 & 5 \\
\hline
\end{tabularx}
\end{table}

\noindent \textit{SL = số lượng kết quả; Top\% = điểm tương đồng cosine cao nhất; Rel. / Rank. = Điểm Relevance / Ranking của LLM Judge (thang 1--5); các dòng nhiễu không tính vào trung bình.}

\subsubsection{Phân tích Hiệu năng Tìm kiếm}

\begin{itemize}
    \item \textbf{NDCG@5 = 0.9288:} Chất lượng xếp hạng đạt mức gần tối ưu. Các cuốn sách phù hợp nhất luôn được ưu tiên đặt lên vị trí đầu trong danh sách kết quả đối với toàn bộ 12 truy vấn ngữ nghĩa.
    \item \textbf{MRR = 0.8750:} Trong 10 trên 12 truy vấn ngữ nghĩa, kết quả tốt nhất xuất hiện ngay ở Vị trí 1 (Rank 1), khẳng định tính chính xác cao ở kết quả đầu tiên (first-hit accuracy).
    \item \textbf{Precision@3 = 0.5556:} Thấp hơn một chút so với mục tiêu 0.60, nguyên nhân chính là do ở các truy vấn S03 (UI/UX) và S11 (Cơ sở dữ liệu), hệ thống không có đủ sách chuyên ngành sát với yêu cầu, dẫn đến việc top 3 kết quả phải trả về các tài liệu mang tính liên quan xa thuộc nhóm Khoa học Máy tính.
    \item \textbf{Noise Rejection = 100.0\%:} Đạt hiệu suất tuyệt đối. Nhờ sử dụng mô hình embedding đa ngôn ngữ chất lượng cao, điểm số cosine phân hóa rất rõ rệt giữa truy vấn hợp lệ ($> 0.55$) và truy vấn nhiễu ($< 0.35$). Do đó, việc thiết lập ngưỡng lọc $\text{threshold} = 0.40$ đã giúp hệ thống tự động chặn đứng toàn bộ các truy vấn không liên quan (trả về 0 kết quả) mà không làm mất đi tài liệu của các truy vấn hợp lệ.
\end{itemize}

